\begin{helper}
Fix one exposure-history cell \(C_h\) from
\(\noderef{FixedSetExposureCellProductLaw}\).  Let \(U_h\) be its terminal
coordinate set, let \(R_h\) be its recorded coordinate set, and let the
coordinates of \(U_h\) be scanned in the fixed terminal order.  Assume the
cell has the cylinder description, the recorded-set identity, the terminal
universe identity, staged-open containment in \(U_h\), the exact product law
for every fixed terminal subset, and the prefix-safety statement for the
terminal order.  Assume also
\[
0\le \varepsilon_1\le 1,\qquad 0\le p\le \frac12,\qquad |I|=k,
\]
and the pointwise lower bound
\[
f(\ell_R,\ell_B)-10\varepsilon_1k^2
\le |\noderef{TwoBiteStagedOpenPairs}(\omega,\varepsilon,I)|
\]
for every regular sample in this projection-size cell.

Assume moreover a branch-indexed adaptive-candidate package.  For every
branch representative \(\sigma\in C_h\), there are terminal sets
\[
E_R(\sigma)\subseteq U_h,\qquad E_B(\sigma)\subseteq U_h
\]
with the following two properties.  If a regular completion \(\omega\in C_h\)
has the same injection and recorded values as \(\sigma\), and agrees with
\(\sigma\) on all blue terminal coordinates, then
\[
\noderef{ClosedPairsBound}(|E_R(\sigma)|,3\varepsilon_1,k),
\]
and every present red staged-open coordinate in an independent final graph
lies in \(E_R(\sigma)\).  Symmetrically, if \(\omega\) agrees with
\(\sigma\) on all red terminal coordinates, then
\[
\noderef{ClosedPairsBound}(|E_B(\sigma)|,3\varepsilon_1,k),
\]
and every present blue staged-open coordinate in an independent final graph
lies in \(E_B(\sigma)\).  Then
\[
\Pr(\mathcal B_I\cap\mathcal R\mid C_h)
\le
\max(1,k^4)\exp\!\left(
8\sqrt{\varepsilon_1}\,pk^2
-p\bigl(f(\ell_R,\ell_B)-10\varepsilon_1k^2\bigr)\right),
\]
where \(\mathcal B_I\) is final-graph independence of \(I\) and
\(\mathcal R\) is the regularity event.
\end{helper}
\begin{proof}
Put
\[
a=f(\ell_R,\ell_B)-10\varepsilon_1k^2.
\]
Since \(C_h\) is contained in the projection-size cell, the assumed open-pair
lower bound says that every regular \(\omega\in C_h\) satisfies
\[
a\le
|\noderef{TwoBiteStagedOpenPairs}(\omega,\varepsilon,I)|. \tag{1}
\]
For such an \(\omega\), let \(P_R(\omega)\) and \(P_B(\omega)\) be the red
and blue terminal staged-open coordinates which are present.  Write their
cardinalities as \(u_R\) and \(u_B\).  The containment
\(\noderef{TwoBiteStagedOpenPairs}(\omega,\varepsilon,I)\subseteq U_h\)
ensures that all these coordinates are read inside the fixed terminal product
set.

Fix \(u_R,u_B\) and first suppose \(u_R\ge u_B\).  The proof pays for the blue
terminal support before conditioning on a complete blue branch; the actual
cylinder scan still follows the fixed terminal order.  More explicitly, fix a
possible blue present support \(B\).  The support \(B\) is a subset of the
projected blue pairs of the \(k\)-set \(I\), so there are at most
\(\binom{k}{2}\) ambient blue coordinates and hence at most
\(\binom{N_B}{u_B}\) such choices with \(N_B\le k^2\).  During the terminal
scan, whenever a coordinate of \(B\) is reached, prefix safety says that the
assertion that it is currently staged-open is already determined from the
history and earlier terminal values; requiring its Bernoulli value to be
present therefore contributes a fresh factor \(p\).  Red coordinates and blue
coordinates outside \(B\) that have not been selected for a required value are
summed out by the exact product law.  Consequently the total probability mass
of all complete blue truth-table branches extending this fixed support \(B\)
is bounded by \(p^{u_B}\Pr(C_h)\), with any additional selected blue absences
contributing their own factors \(1-p\) when they are selected later.  This is
the point at which the first-color support probability is paid; it is not
recovered after conditioning on a complete branch.

Now refine this first-color support event into complete blue terminal
branches.  For a nonempty branch choose a representative
\(\sigma\in C_h\) with exactly that blue terminal truth table, and use the
branch-indexed hypothesis to obtain \(E_R(\sigma)\subseteq U_h\).  For every
regular independent completion extending this complete blue branch, all
present red staged-open coordinates lie in \(E_R(\sigma)\), and
\[
\noderef{ClosedPairsBound}(|E_R(\sigma)|,3\varepsilon_1,k)
\]
gives
\[
|E_R(\sigma)|\le 3\varepsilon_1k^2. \tag{2}
\]
Thus, after the complete blue branch is fixed, there are at most
\(\binom{N_R}{u_R}\) choices for the red present support \(R\), with
\(N_R\le 3\varepsilon_1k^2\).

Fix such \(B\), a complete blue branch extending \(B\), and a red support
\(R\subseteq E_R(\sigma)\) of size \(u_R\).  For any successful completion
with these data, (1) and the definitions of \(u_R,u_B\) imply that at least
\[
a-u_R-u_B
\]
staged-open terminal coordinates outside \(R\cup B\) are absent.  If this
number is nonpositive, no absence condition is needed for an upper bound.  If
it is positive, let \(r=\lceil a-u_R-u_B\rceil\).  The number of absent
staged-open coordinates outside \(R\cup B\) is an integer at least
\(a-u_R-u_B\), hence at least \(r\).

To make the absence condition canonical, scan the fixed terminal order and
select the first \(r\) coordinates which the prefix-measurable staged-open
test declares to be outside \(R\cup B\).  Prefix safety is used here
coordinate by coordinate: membership in the staged-open set, touchedness,
same-color closedness, and recordedness at the coordinate currently being
tested depend only on the history and previously exposed terminal values, not
on the coordinate's own Bernoulli value or on later terminal values.  Therefore
each selected coordinate is still a fresh Bernoulli trial when it is required
to be absent.  If a selected absent coordinate has blue color, its
\((1-p)\)-factor is part of the summed blue-branch mass described above; if it
has red color, it is charged in the branch-relative red exposure.  In either
case the combined cylinder specified by \(B\), \(R\), and the selected absent
coordinates has probability at most
\[
p^{u_B}p^{u_R}(1-p)^r\Pr(C_h)
\]
before normalization by \(\Pr(C_h)\).  Unmentioned terminal coordinates are
summed out by the exact product law, and the additional regularity condition
only restricts the event further.

Thus, for this fixed count pair in the case \(u_R\ge u_B\), the conditional
probability is bounded by
\[
\binom{N_R}{u_R}p^{u_R}
\binom{N_B}{u_B}p^{u_B}
(1-p)^{\max(0,a-u_R-u_B)}
\]
with \(N_R\le3\varepsilon_1k^2\) and \(N_B\le k^2\).  The complete blue
branches are used only to choose the branch-dependent exceptional set
\(E_R(\sigma)\).  Their probabilities remain the averaging weights, exactly
as in \(\noderef{FixedSetHistoryCellBranchAveraging}\), while the separate
first-color support calculation above supplies the \(p^{u_B}\) factor.  Hence
there is no hidden factor for the number of blue branches and no fresh
first-color factor is charged after conditioning on a complete branch.

Applying
\(\noderef{FixedSetHistoryCellBinomialExponentialEstimate}\) gives
\[
\Pr(\mathcal B_I\cap\mathcal R
\cap\{|P_R|=u_R,\ |P_B|=u_B\}\mid C_h)
\le
\exp\!\left(
8\sqrt{\varepsilon_1}\,pk^2-pa\right). \tag{3}
\]
The helper handles \(p=0\) before introducing \(2p\), and it treats the case
\(a-u_R-u_B<0\) with absence exponent \(0\) rather than manipulating a
negative absence exponent.

When \(u_B\ge u_R\), the same two-stage argument is run with the colors
interchanged.  One first pays for the red support mass in the fixed terminal
order, then refines into complete red terminal branches in order to obtain
\(E_B(\sigma)\), and finally charges the blue support and the canonical absent
coordinates.  The same estimate (3) holds for every fixed ordered pair
\((u_R,u_B)\).

It remains only to sum over possible count pairs.  The projected pairs of a
\(k\)-element set are unordered pairs represented by increasing indices, so
each color has at most \(\binom{k}{2}\) staged-open terminal coordinates
available.  Hence the possible values of \(u_R\) are contained in
\[
\{0,1,\ldots,\binom{k}{2}\},
\]
and similarly for \(u_B\).  If \(k=0\), this gives exactly one ordered count
pair, and the factor \(\max(1,k^4)\) covers it.  If \(0<k\), then
\[
\binom{k}{2}+1\le k^2:
\]
for \(k=1\) both sides are \(1\), while for \(k\ge2\),
\(\binom{k}{2}+1\le \binom{k}{2}+k\le k^2\).  Thus there are at most \(k^2\)
possible values for each color and at most \(k^4\) ordered count pairs.  In
all cases, the number of ordered count pairs is at most \(\max(1,k^4)\).
Finite subadditivity of \(\noderef{TwoBiteConditionalProbability}\) on these
disjoint count events, together with the uniform bound (3), gives
\[
\Pr(\mathcal B_I\cap\mathcal R\mid C_h)
\le
\max(1,k^4)\exp\!\left(
8\sqrt{\varepsilon_1}\,pk^2-pa\right).
\]
Substituting the definition of \(a\) is exactly the claimed bound.
\end{proof}
